% =====================================================================
% Stage-2 conditional re-forecast ("Frozen 2"). Frozen methods section.
% The method is pre-registered (paper/STAGE2_PREREG.md); the empirical
% tables in sec:stage2-results are filled after the 2026 knockout phase.
% Code is hash-pinned in the pre-registration (branch ko-pool-repick).
% =====================================================================

\section{A second freeze: the Stage-2 conditional re-forecast}\label{sec:stage2}

The pre-registered forecast of Section~\ref{sec:evolution} (``Frozen~1'') commits to
a full bracket before the tournament and never changes. The pool, however,
permits each entrant to re-pick the entire knockout bracket once the group stage
concludes, and so it offers a second decision point that the pre-tournament entry
cannot use: the realized group results, the teams that actually advanced, and the
draw they actually face. We leverage this second decision point to issue a
\emph{Stage-2 conditional re-forecast} (``Frozen~2''), a knockout entry that
conditions on everything known at the group$\rightarrow$knockout boundary and is
optimized for expected pool points on the bracket that actually formed. Frozen~2
is a distinct object from Frozen~1. It uses look-ahead relative to the
pre-tournament entry by construction, so it is not pre-tournament
pre-registration; it is a second freeze, of the \emph{method}, fixed before any
knockout result is observed (Section~\ref{sec:stage2-prereg}). Frozen~1 and the
paper's pre-registered forecast are unchanged.

\subsection{The scoring rule and the re-pick problem}\label{sec:stage2-rule}

The pool scores a knockout match in two parts. The first is the
\emph{90-minute regulation} scoreline, graded on the same three-tier rule as the
group stage: three points for the exact score, two for the correct result and
goal difference, and one for the correct result alone. The second is a separate
point for naming the team that advances. Extra-time and penalty \emph{scores} are
not graded; they enter only through the advancing team. A bracket is submitted
forward, from the Round of~32 to the final, and a later-round pick is graded only
if the projected matchup actually occurs, because the entry commits to specific
teams meeting in a specific slot. The objective is to maximize the expected total
of these per-match points over the realized bracket.

This structure makes the re-pick a coupled problem rather than a sequence of
independent guesses. The advancer named in one match determines which teams
populate the next slot, and a later pick earns nothing unless both of its
projected teams arrive there, so the value of a deep pick is the product of its
per-match value and the probability that its matchup materializes. We treat the
two pieces separately: a per-match forecast (Section~\ref{sec:stage2-match}) and
the bracket that assembles them (Section~\ref{sec:stage2-dp}).

\subsection{Ratings and the per-match forecast}\label{sec:stage2-match}

The Stage-2 forecast conditions on post-group information through the ratings.
We rate each team by its Track~B effective Elo, which combines the live Elo, the
current injury and lineup adjustments, and the learned in-tournament drift of
Appendix~\ref{app:prophet}, and carries no separate host bonus because the host
advantage is already absorbed into the live ratings. For a matchup we map the
rating difference to a pair of expected goals through the same logistic
win-expectancy used throughout the paper, and model the 90-minute score as
independent Poisson draws on those means.

Two quantities follow from this match model. The first is the expected-value
scoreline under the three-tier rule, which is the group-stage object of
Equation~\eqref{eq:ev} applied to the 90-minute distribution. The second is the
probability that a team advances, which the 90-minute distribution alone does not
give, because a drawn regulation is decided later. We compute it as a cascade:
the team advances if it wins in 90 minutes, or the match is level after 90 and it
wins in extra time, or the match is level after extra time and it wins the
shootout. We model extra time as a Poisson period with the means scaled to a
third of the regulation rate, the 30 added minutes against 90, and we set the
shootout win probability to one half. The coin flip is deliberate: in
Section~\ref{sec:audit} we find that a team's prior shootout record does not
predict the next shootout, so there is no signal to exploit and a tilt would only
add noise.

The per-match pick maximizes the sum of the two pieces, the scoreline's
expected three-tier points plus the probability that the named advancer
advances. The two choices are coupled by the rule itself, because a decisive
90-minute prediction forces the advancer to the predicted winner while a
predicted draw leaves the advancer free. The optimizer therefore weighs a
decisive scoreline against a draw that banks the advancing point on the
likelier team, and on the most even matchups it prefers the draw, consistent
with the group-stage result that draws become competitive when the result is
close.

\subsection{The bracket as a dynamic program}\label{sec:stage2-dp}

Because a deep pick is graded only if its matchup occurs, the expected value of
the whole entry is a sum over slots of the per-match value weighted by the
probability that the projected pair occupies the slot. That occupancy
probability has a useful structure. The two matches feeding a slot draw on
disjoint halves of the bracket, so the team emerging from one feed is
independent of the team emerging from the other, and the occupancy weight
factors into the product of two one-sided reaching probabilities. Those
probabilities are properties of the realized bracket alone and not of the
entry, so we estimate them once by Monte~Carlo and then solve for the entry
exactly.

We compute the globally optimal internally-consistent bracket by dynamic
programming over the bracket tree. Working from the round of~32 upward, for each
match and each team the entry could project out of it, we store the best total
expected points obtainable in that match's subtree, and we combine subtrees at
their parent. The greedy alternative, which simply advances the stronger team in
each match, is not optimal, because advancing the stronger team maximizes the
local advancing point but ignores how that choice changes both the reachability
and the matchup value of later rounds. On the realized 2026 bracket we report the
dynamic program's improvement over the greedy entry, and the handful of
coin-flip slots that drive it, with the post-knockout results
(Section~\ref{sec:stage2-results}).\footnote{On the predicted pre-knockout
bracket the improvement is small, on the order of a tenth of a point, and
changes a single advancer. We read this as a located refinement rather than a
large gain, because most knockout matchups are lopsided enough that the greedy
and optimal advancers coincide; the value of the re-pick is in re-conditioning on
the realized draw, not in the optimizer's last increment.}

\subsection{Independence and robustness}\label{sec:stage2-robust}

We keep the independent-Poisson match model rather than adding a Dixon--Coles
low-score correlation, and we verify that this does not bind on the choice we
care about. Across the bracket's 32 knockout matchups the correction changes the
expected-value scoreline in none of them for correlation parameters in
$\{0.05, 0.10, 0.15\}$, because the optimal picks sit at scorelines the
correction reweights too little to overturn, consistent with the group-stage
audit of Section~\ref{sec:audit}. The extra-time scaling and the shootout coin
flip enter only the residual draw-decided mass and move no advancer, so the entry
is robust to those assumptions as well.

\subsection{Pre-registration}\label{sec:stage2-prereg}

We froze the Stage-2 method before the first knockout match, and recorded the
freeze publicly so that the re-pick reads as a second forecast and not a
hindsight fit. The record fixes the scoring rule, the cascade, the dynamic
program, the Track~B ratings, and the hyperparameters that were tuned on the
2018 and 2022 tournaments and not on any 2026 knockout outcome, and it pins the
implementation by hash. It does not change Frozen~1, the pre-registered living
forecast, or the locked group-stage picks. Holding the method fixed across the
boundary is what lets the post-knockout comparison in
Section~\ref{sec:stage2-results} be read as the value of re-conditioning rather
than the value of re-optimizing the method with the answers in hand.

\subsection{What we report after the knockout phase}\label{sec:stage2-results}

% STUB --- the empirical tables below are filled after the 2026 knockout phase.
The Stage-2 forecast is evaluated only after the knockout phase, on three
questions. First, whether the realized pool points of the Frozen~2 entry exceed
those of the greedy, naive-favorite, and Frozen~1 knockout brackets. Second, the
value of re-conditioning, the Frozen~2 minus Frozen~1 realized knockout points,
which isolates the gain from conditioning on the group stage holding the method
fixed. Third, calibration, whether the entry's realized points track its expected
points. We also report a full 2018 and 2022 backtest with a 32-team bracket
adapter, post-group learned Elo, and true 90-minute scores; a lean version,
holding the favorite-advancing bracket fixed and varying only the scoreline rule,
scores the expected-value entry at 33 points against 32 for naive favorite-1-0
picks across the two tournaments, which is positive but small and consistent with
the group-stage finding that the expected-value scoreline narrowly beats the
naive baseline. The lean check is superseded by the full backtest here.

% \input{tables/stage2_points}    % EV-optimal vs greedy / naive / Frozen 1
% \input{tables/stage2_recond}    % value of re-conditioning (F2 - F1)
% \input{tables/stage2_dp}        % DP vs greedy gain, by slot
